\section{Governance, Roles, and Sustainability}\label{sec:governance}

Governance is intentionally lean: few bodies, clear roles, and short decision paths.

\subsection{Roles and responsibilities}
\begin{longtable}{@{}p{3.6cm}p{11.4cm}@{}}
\toprule
Role & Responsibilities \\
\midrule
\endhead
Policy Owner and governing body: Technical Working Group for AI Policy and Governance (TWG), led by Dr.\ Magdalene C.\ Unajan & Owns this framework and is accountable for its content, interpretation, review, and amendment. Also performs the functions of a steering body: reviews Tier~C tools and exceptions, handles interpretive questions, and conducts the annual review. The TWG may consult the Data Protection Officer, the Information Technology Office, and faculty and student representatives as needed. \\
AI Support Desk (function designated by the TWG) & Single point of contact; maintains the tool guidance list, template bank, and shared practice repository; collects feedback and incident reports. \\
Data Protection Officer (DPO) & Advises on privacy impact for Tier~B and C tools and handles personal-data incidents under the Data Privacy Act. \\
College and Department Heads & Promote adoption, release time for professional development where feasible, and handle first-line questions. They do not add blanket restrictions beyond this framework without TWG consent. \\
Faculty & Decide AI-use levels for their courses (Section~\ref{sec:integrity}), disclose AI use where required, and model responsible practice. \\
Students & Follow the course AI-use declaration, disclose AI use as required, and verify AI outputs they rely on. \\
Information Technology Office & Maintains the approved tool list, access arrangements, and security controls. \\
Student Affairs / Integrity Office & Administers the proportionate process for suspected misuse (Section~\ref{sec:integrity}). \\
\bottomrule
\end{longtable}

\subsection{Decision principles}
\begin{itemize}
  \item Course-level AI-use decisions rest with the faculty member within the framework limits.
  \item Tier~A uses need no approval; Tier~B requires use of listed tools;   Tier~C requires TWG review (Section~\ref{sec:tools}).
    \item The TWG publishes its review turnaround targets (Section~\ref{sec:tools}); working days exclude holidays, special non-working days, and days when work is suspended.
\end{itemize}

\subsection{Document control}
\begin{itemize}
  \item \textbf{Policy owner:} Technical Working Group for AI Policy and Governance, led by Dr.\ Magdalene C.\ Unajan.
  \item The framework carries a version number (major.minor), a change log, and an effective date.
  \item Minor updates (clarifications, tool list changes, template updates) are approved by the TWG and logged. Major changes follow the amendment procedure (Section~\ref{sec:amendment}).
  \item Supporting materials (tool list, templates, quick guides) are living documents kept outside the policy text so frequent updates do not require amendment.
\end{itemize}

\subsection{Review cycle and feedback loop}
\begin{enumerate}
  \item \textbf{Continuous feedback:} a simple form and the Support Desk collect feedback and incidents year-round.
  \item \textbf{Semestral check-in:}   the Support Desk reports brief metrics and emerging issues to the TWG.
    \item \textbf{Annual review:} the TWG reviews metrics, feedback, and technological and regulatory developments, and recommends updates.
  \item \textbf{Triggered review:} a significant legal change, CHED issuance, serious incident, or major technology change prompts an out-of-cycle review.
\end{enumerate}
